\chapter{Methodology}
\label{ch:methodology}

This chapter describes, in full reproducible detail, the data,
algorithms, and protocols that underpin the system described in
Chapter~\ref{ch:architecture}. Section~\ref{sec:m-wesad} presents
the WESAD dataset and the rationale for its use. Section~\ref{sec:m-pre}
covers signal pre-processing (resampling, IBI extraction, windowing).
Section~\ref{sec:m-labels} specifies the label mapping from the raw
five-class WESAD protocol to our four ordinal stress classes.
Sections~\ref{sec:m-arch} and \ref{sec:m-train} describe the neural
network architecture and the training procedure respectively.
Section~\ref{sec:m-tflite} details the TensorFlow Lite conversion
pipeline, and Sections~\ref{sec:m-ble} and \ref{sec:m-fsm} specify
the BLE protocol and the Arduino finite-state machine.

\section{WESAD Dataset}
\label{sec:m-wesad}

The WESAD dataset \citep{schmidt2018wesad} contains synchronised
recordings from fifteen healthy adult subjects (11~male, 4~female,
mean age $27.5 \pm 2.4$ years) wearing two devices simultaneously: a
RespiBAN Professional chest band and an Empatica E4 wrist wearable.
Subjects are indexed \texttt{S2} through \texttt{S17}; subject
\texttt{S12} is missing from the public release because of
instrumentation failure during the recording session, leaving
fifteen usable subjects. Each subject follows a fixed approximately
one-hour protocol consisting of four affective conditions -- baseline,
amusement, stress, and meditation -- separated by brief transition
intervals. The stress condition is implemented as a variant of the
Trier Social Stress Test in which subjects are required to deliver
a five-minute impromptu speech followed by five minutes of mental
arithmetic in front of an evaluation panel.

Because our deployment target is a wrist wearable, we use
\emph{only} the Empatica E4 wrist channels and we ignore the chest
channels entirely. Table~\ref{tab:wesad-sensors} lists the four
wrist channels and their native sampling rates.

\begin{table}[H]
\centering
\caption{Empatica E4 wrist sensors in the WESAD dataset
\citep{schmidt2018wesad} used for this project.}
\label{tab:wesad-sensors}
\small
\begin{tabular}{lcl}
\toprule
\textbf{Sensor} & \textbf{Sampling rate (Hz)} & \textbf{Description} \\
\midrule
BVP (PPG) & 64 & Blood volume pulse, green-light photoplethysmography \\
EDA & 4 & Electrodermal activity, wrist electrodes \\
TEMP & 4 & Skin temperature at the wrist, NTC thermistor \\
ACC (3-axis) & 32 & Triaxial acceleration, $\pm 2$g \\
Labels & 700 & Condition label per sample (0--4) \\
\bottomrule
\end{tabular}
\end{table}

All four channels are available for all fifteen subjects, and the
label channel is synchronised at 700~Hz and downsampled to the
channel-specific rates as needed during pre-processing.

\section{Signal Pre-processing}
\label{sec:m-pre}

Our pre-processing pipeline has three sequential stages: per-channel
resampling to a common 4~Hz grid, inter-beat-interval (IBI) extraction
from BVP, and sliding-window segmentation.

\textbf{Resampling.} The four channels in Table~\ref{tab:wesad-sensors}
sample at different rates, which is incompatible with a single
tensor representation. We resample each channel to a common
\textbf{4~Hz} grid -- the same rate as EDA and TEMP natively -- for
two reasons. First, stress-relevant dynamics operate on timescales
of seconds to tens of seconds, and a 4~Hz rate preserves more than
enough bandwidth to capture them. Second, a 4~Hz common rate keeps
the window tensor small ($30 \times 5 = 150$ floats) and therefore
compatible with the memory budget of a Wear~OS device. The BVP channel
is downsampled from 64~Hz to 4~Hz using SciPy's polyphase
\texttt{resample\_poly} function (up=1, down=16), which applies an
anti-aliasing low-pass filter designed for exactly this ratio. The
ACC channel is downsampled from 32~Hz to 4~Hz by first computing
the per-sample magnitude $\sqrt{a_x^2 + a_y^2 + a_z^2}$ and then
polyphase-resampling by a factor of 8. EDA and TEMP are already at
4~Hz and are used directly.

\textbf{IBI extraction.} A separate IBI stream is derived from the
native 64~Hz BVP before it is downsampled. We use SciPy's
\texttt{find\_peaks} with a minimum prominence of 0.1 times the
channel standard deviation and a minimum distance equal to 0.4
seconds' worth of samples, which corresponds to a physiologically
plausible maximum heart rate of 150~bpm. Successive peak-to-peak
intervals yield an irregular IBI time series which is then
forward-filled onto the 4~Hz master grid: at each 4~Hz step we emit
the most recent IBI value. Forward-filling is appropriate here
because IBI is a ``most recent beat'' quantity whose validity
persists until a new beat is detected.

\textbf{Windowing.} The five-channel 4~Hz signal (BVP, IBI, EDA, TEMP,
ACC magnitude) is segmented into overlapping windows of $30$ samples
(7.5 seconds of wall-clock time) with a step of $15$ samples (50\%
overlap). Each window becomes one training example, and its
majority-vote label across its 30 samples becomes the supervisory
target. Windows whose majority label is ``transient''
(condition codes not belonging to any of the four target conditions)
are discarded. This yields approximately 5{,}000 usable windows
aggregated across all fifteen subjects, which are then split 80/20
into train and test partitions before normalisation.

\section{Label Mapping}
\label{sec:m-labels}

WESAD uses a five-valued label channel in which 0 denotes an
invalid or transient sample, 1 the baseline condition, 2 the
stress condition, 3 the amusement condition, and 4 the meditation
condition. For the purposes of this project we collapse the two
condition codes that are not associated with negative arousal
(baseline and meditation) into a single Calm class, we treat
amusement as Mild positive arousal, and we \emph{split} the stress
condition into two ordinal classes based on physiological severity.
The full mapping is shown in Table~\ref{tab:label-map}.

\begin{table}[H]
\centering
\caption{Mapping from WESAD condition codes to the four ordinal
stress classes used in this project.}
\label{tab:label-map}
\small
\begin{tabular}{clll}
\toprule
\textbf{WESAD code} & \textbf{WESAD condition} & \textbf{Our class} &
\textbf{Severity} \\
\midrule
0 & Transient / invalid & (discarded) & -- \\
1 & Baseline            & Calm        & 0 \\
4 & Meditation          & Calm        & 0 \\
3 & Amusement           & Mild        & 1 \\
2 & Stress (normal physio) & Moderate & 2 \\
2 & Stress (EDA $>$ 75\%ile \& IBI $<$ 25\%ile) & Critical & 3 \\
\bottomrule
\end{tabular}
\end{table}

The stress-split rule is as follows. We compute, \emph{per subject},
the 75th percentile of EDA across all stress-labelled samples of
that subject (call it $\theta_\text{EDA}^s$) and the 25th percentile
of IBI across all stress-labelled samples of that subject (call it
$\theta_\text{IBI}^s$). A stress-labelled sample belonging to subject
$s$ is classified as \emph{Critical} if and only if its EDA exceeds
$\theta_\text{EDA}^s$ and simultaneously its IBI falls below
$\theta_\text{IBI}^s$. All other stress-labelled samples are
classified as \emph{Moderate}. The per-subject quantile design is
deliberate: it corrects for between-subject differences in baseline
EDA amplitude and resting heart rate, which are known to vary by
more than a factor of two across healthy adults, and it ensures
that ``Critical'' denotes a within-subject physiological extreme
rather than an absolute one.

\section{Neural Network Architecture}
\label{sec:m-arch}

The neural network is a compact bidirectional recurrent classifier
with an additive attention pooling layer. Figure~\ref{fig:model-arch}
shows a schematic; Table~\ref{tab:layers} lists every layer with its
output shape and approximate parameter count.

\begin{figure}[H]
\centering
\includegraphics[width=\linewidth]{model_architecture}
\caption{Neural network architecture for the four-class driver stress
classifier. The input is a 30-step, 5-feature window. A bidirectional
LSTM \citep{graves2005bilstm} produces a 128-dimensional hidden state
per step; an additive attention layer \citep{vaswani2017attention}
collapses the sequence to a single 128-dimensional context vector; a
two-layer feed-forward head maps to a four-class softmax.}
\label{fig:model-arch}
\end{figure}

\begin{table}[H]
\centering
\caption{Per-layer breakdown of the stress classifier.}
\label{tab:layers}
\small
\begin{tabular}{lll}
\toprule
\textbf{Layer} & \textbf{Output shape} & \textbf{Parameters} \\
\midrule
Input                     & $[1, 30, 5]$   & 0      \\
BiLSTM(64, return\_seq=True) & $[1, 30, 128]$ & 35{,}840 \\
Dropout(0.3)              & $[1, 30, 128]$ & 0      \\
Attention: Dense(1, tanh) & $[1, 30, 1]$   & 129    \\
Attention: Softmax (axis=timesteps) & $[1, 30, 1]$ & 0 \\
Temporal weighted sum     & $[1, 128]$     & 0      \\
Dense(64, relu)           & $[1, 64]$      & 8{,}256 \\
Dropout(0.3)              & $[1, 64]$      & 0      \\
Dense(32, relu)           & $[1, 32]$      & 2{,}080 \\
Dense(4, softmax)         & $[1, 4]$       & 132    \\
\midrule
\textbf{Total}            &                & $\approx 160$K \\
\bottomrule
\end{tabular}
\end{table}

The choice of a bidirectional LSTM backbone follows
\citet{graves2005bilstm}: bidirectionality allows every hidden state
to incorporate both past and future context within the 7.5-second
window, which is particularly useful because autonomic responses
lag their triggers by several seconds. The additive attention layer
is inspired by the scaled-dot-product attention of
\citet{vaswani2017attention}, but in a minimal single-head form
that adds only 129 parameters. It learns a scalar importance weight
per timestep, and the softmax across timesteps enforces that these
weights sum to one, producing a convex combination of the 30 hidden
states. Empirically the attention distribution concentrates on the
timesteps where EDA peaks or IBI shortens, which matches physiological
intuition.

\section{Training Procedure}
\label{sec:m-train}

The model is trained end-to-end using stochastic gradient descent
with the Adam optimiser. All hyperparameters are summarised in
Table~\ref{tab:hparams}.

\begin{table}[H]
\centering
\caption{Training hyperparameters.}
\label{tab:hparams}
\small
\begin{tabular}{ll}
\toprule
\textbf{Hyperparameter} & \textbf{Value} \\
\midrule
Maximum epochs       & 60 \\
Batch size           & 32 \\
Optimiser            & Adam \\
Initial learning rate & $10^{-3}$ \\
Loss                 & Categorical cross-entropy \\
EarlyStopping patience & 10 (monitor: val loss) \\
ReduceLROnPlateau patience & 5 \\
ReduceLROnPlateau factor & 0.5 \\
Train/test split     & 80 / 20, subject-stratified \\
Normalisation        & z-score per feature, fit on train only \\
\bottomrule
\end{tabular}
\end{table}

A critical implementation detail is \emph{normalisation-after-split}:
we compute the per-feature mean and standard deviation using
\emph{only} the training partition and then apply those statistics
unchanged to the test partition. This prevents information leakage
from test to train, which would otherwise inflate test accuracy
by roughly one to two percentage points. The learned normalisation
constants are also exported alongside the TensorFlow Lite model so
that the Wear OS application can apply the same transformation
at inference time.

EarlyStopping with patience 10 and ReduceLROnPlateau with patience 5
together keep the model from overfitting the small dataset: in
practice the learning rate is halved once or twice during training,
and training terminates between epoch 30 and epoch 45 depending on
the random seed, well before the sixty-epoch maximum.

\section{TensorFlow Lite Conversion}
\label{sec:m-tflite}

Deployment on Wear~OS requires the trained Keras model to be
converted to the TensorFlow Lite \citep{tflite2023} FlatBuffer
format. The conversion follows three steps.

\textbf{Step 1: SavedModel export.} The trained Keras model is
exported using \texttt{model.save(path, save\_format='tf')},
producing a SavedModel directory that preserves the full computational
graph including the bidirectional recurrent structure.

\textbf{Step 2: \texttt{TFLiteConverter.from\_saved\_model}.} We
instantiate the converter from the SavedModel directory and enable
two flags:
\begin{itemize}
  \item \texttt{optimizations = [tf.lite.Optimize.DEFAULT]}, which
  enables dynamic-range post-training quantisation and, in this
  configuration, reduces the model footprint by roughly a factor
  of four;
  \item \texttt{target\_spec.supported\_ops =
  [tf.lite.OpsSet.TFLITE\_BUILTINS, tf.lite.OpsSet.SELECT\_TF\_OPS]}.
\end{itemize}

The \texttt{SELECT\_TF\_OPS} flag is strictly required because the
bidirectional LSTM unrolls internally into \texttt{TensorListReserve},
\texttt{TensorListFromTensor}, and \texttt{TensorListStack} operations
that are not part of the core TFLite builtin set. The
\texttt{SELECT\_TF\_OPS} path includes a small runtime shim that
executes these operations at inference time; the shim adds roughly
30~KB to the on-device TensorFlow Lite runtime footprint, which is
a negligible cost on any Wear~OS device and is a one-time cost rather
than per-model.

\textbf{Step 3: Write the FlatBuffer.} The converter emits a
\texttt{bac\_model.tflite}-equivalent \texttt{stress\_model.tflite}
FlatBuffer of approximately \textbf{70~KB}, well under the 80~KB
target and within the 22~MB application budget for the Wear~OS APK.
The same conversion pipeline is used unchanged for the AlcoWatch
BAC model and was validated in Chapter~\ref{ch:architecture}'s
cross-reference against \citet{acn1408}.

\section{BLE Protocol}
\label{sec:m-ble}

The wearable exposes a single GATT service containing one notifiable
characteristic. Table~\ref{tab:ble-packet} defines the twelve-byte
characteristic payload.

\begin{table}[H]
\centering
\caption{Twelve-byte BLE stress characteristic payload. Service UUID
\texttt{ABCD1234-1234-5678-1234-56789abcdef0}; characteristic UUID
\texttt{ABCD1235-1234-5678-1234-56789abcdef0}.}
\label{tab:ble-packet}
\small
\begin{tabular}{clll}
\toprule
\textbf{Offset} & \textbf{Size} & \textbf{Field} & \textbf{Encoding} \\
\midrule
0  & 8 bytes & Timestamp      & int64, little-endian, ms since epoch \\
8  & 1 byte  & Stress level   & uint8, 0=Calm, 1=Mild, 2=Moderate, 3=Critical \\
9  & 1 byte  & Confidence     & uint8, confidence $\times 100$, range 0--100 \\
10 & 2 bytes & Reserved       & zero-initialised, reserved for future use \\
\bottomrule
\end{tabular}
\end{table}

The packet is notified to the subscribed central once every thirty
seconds. Because BLE notifications are unacknowledged at the
attribute-protocol level, a reserved field and an explicit timestamp
are included so that the central can detect stale or out-of-order
packets and a future revision of the protocol can carry additional
signals (for example, a per-class probability distribution) without
breaking existing clients.

\section{Arduino Cabin Finite-State Machine}
\label{sec:m-fsm}

The Arduino firmware implements the five-state machine summarised
in Table~\ref{tab:fsm}.

\begin{table}[H]
\centering
\caption{Arduino cabin finite-state machine. All transitions are driven
by incoming BLE notifications; the \texttt{WAITING} state is additionally
reachable from any other state after a 60-second timeout.}
\label{tab:fsm}
\small
\begin{tabular}{lllll}
\toprule
\textbf{State} & \textbf{Condition} & \textbf{LED (RGB)} &
\textbf{Fan (PWM \%)} & \textbf{Audio} \\
\midrule
WAITING  & boot OR no packet $\ge 60$s & yellow blink & 0    & off \\
CALM     & level $= 0$                 & warm white   & 30\% & off \\
MILD     & level $= 1$                 & neutral white & 50\% & off \\
MODERATE & level $= 2$                 & blue          & 75\% & soft music \\
CRITICAL & level $= 3$                 & red blinking  & 100\% & voice alert \\
\bottomrule
\end{tabular}
\end{table}

Two aspects of the FSM deserve emphasis. First, the \texttt{WAITING}
state is the \emph{safe default}: at boot, when the wearable is out
of range, when the watch battery has died, or when the driver has
removed the watch, the cabin enters a visually neutral low-intervention
state rather than committing to an unjustified response. The sixty-second
BLE timeout was chosen to be twice the nominal thirty-second emission
cadence, which gives the firmware headroom for one missed notification
without dropping the cabin into \texttt{WAITING}. Second, the severity
ordering of the four active states -- from warm white (Calm) through
neutral white (Mild) and calming blue (Moderate) to blinking red
(Critical) -- is \emph{monotonic} in perceived intensity. This gives
the driver an unambiguous cue that the cabin is tracking a progressive
physiological quantity, rather than a categorical distinction without
ordinal meaning.
